% Common preamble for all lab reports
\documentclass[a4paper, 12pt]{article}

% Language and encoding
\usepackage[T2A]{fontenc}
\usepackage[utf8]{inputenc}
\usepackage[ukrainian]{babel}

% Typography
\usepackage{microtype}
\usepackage{lmodern}

% Page geometry
\usepackage[margin=25mm, headheight=14pt]{geometry}

% Math
\usepackage{amsmath, amssymb}

% Graphics and figures
\usepackage{graphicx}
\usepackage{float}
\usepackage{subcaption}
\usepackage[justification=centering]{caption}

% Tables
\usepackage{tabularx}
\usepackage{booktabs}
\usepackage{array}

% Colors and hyperlinks
\usepackage{xcolor}
\usepackage[colorlinks=true, linkcolor=blue!60!black, urlcolor=blue!60!black, citecolor=blue!60!black]{hyperref}

% Lists
\usepackage{enumitem}
\setlist{nosep, leftmargin=*}

% Section formatting
\usepackage{titlesec}
\titleformat{\section}{\large\bfseries}{\thesection.}{0.5em}{}
\titleformat{\subsection}{\normalsize\bfseries}{\thesubsection.}{0.5em}{}

% Header/footer
\usepackage{fancyhdr}
\pagestyle{fancy}
\fancyhf{}
\fancyhead[L]{\small\leftmark}
\fancyhead[R]{\small\thepage}
\renewcommand{\headrulewidth}{0.4pt}
\fancypagestyle{plain}{\fancyhf{}\fancyfoot[C]{\thepage}\renewcommand{\headrulewidth}{0pt}}

% Placeholder image command
\newcommand{\placeholder}[2][0.8\textwidth]{%
  \begin{figure}[H]
    \centering
    \fbox{\parbox{#1}{\vspace{4cm}\centering\textit{#2}\vspace{4cm}}}
    \caption{#2}
  \end{figure}
}

% Placeholder for inline result values
\newcommand{\result}[1]{\textbf{\textcolor{red!70!black}{[#1]}}}

\usepackage{../common/labstyle}

\labnum{4}
\labtitle{Ультразвук}

\begin{document}
\maketitle
\tableofcontents
\newpage

% ===================================================================
\section{Вихiднi данi}
% ===================================================================

\begin{table}[H]
  \centering
  \caption{Вихiднi данi (варiант~13)}
  \begin{tabularx}{0.75\textwidth}{Xr}
    \toprule
    \textbf{Параметр} & \textbf{Значення} \\
    \midrule
    Частота випромiнювача ($f$)     & 23~кГц \\
    Амплiтуда коливань ($A$)        & 1{,}3~мкм \\
    Об'єм ванночки                  & до 5~л \\
    Рiдина                          & CompressibleWater \\
    Тип розрахунку                  & Transient \\
    \bottomrule
  \end{tabularx}
\end{table}

\subsection{Параметри хвилi}

Довжина хвилi:
$\lambda = c/f = 1480/23000 \approx 64{,}3$~мм.

Розмiр елемента сiтки: $h_{\max} \leq \lambda/10 \approx 6{,}4$~мм (використано 3~мм $\approx \lambda/21$).

Часовий крок: $\Delta t = 0{,}5$~мкс ($\leq 1/(20f) \approx 2{,}17$~мкс).

Гармонiчний рух випромiнювача:
\begin{equation}
  z(t) = A\,\sin(2\pi f\, t) = 1{,}3 \times 10^{-3}\,\sin(2\pi \times 23000 \times t)~\text{[мм]}.
\end{equation}

% ===================================================================
\section{Послiдовнiсть виконання}
% ===================================================================

Робота виконана через ANSYS Workbench з солвером CFX (Transient).

\begin{lstlisting}[style=scheme]
1. Geometry: CadQuery (geometry/lab4.py)
   |-- Stadium bath: 220x120x150mm (capsule shape)
   |-- Bottom fillet R=10mm
   |-- Emitter recess: D=40mm, depth=5mm at center
   |-- Export STEP -> geometry/step/lab4.step
2. ANSYS Workbench: CreateSystem("Fluid Flow (CFX)")
3. Import lab4.step -> SpaceClaim (1 body)
4. Enter Meshing (Python via SendCommand)
   |-- NS "Emitter" -> face at Z=5mm, centroid at
   |   origin, area ~1257mm^2 (circle D=40mm)
   |-- NS "Walls"   -> all remaining faces
   |-- Mesh 3mm (lambda/21) -> 200418 nodes,
   |   1095393 elements
5. CFX-Pre (CCL import)
   |-- MATERIAL: CompressibleWater
   |     EOS = Linear Elastic, c=1480 m/s,
   |     rho=998.2 kg/m^3
   |-- EXPRESSIONS: freq=23kHz, amp=1.3um,
   |     disp_z = amp*sin(2*pi*freq*t)
   |-- FluidDomain: CompressibleWater, Laminar
   |     Mesh Deformation = Regions of Motion Specified
   |     Emitter: Specified Displacement Z=disp_z
   |     Walls: Stationary, No Slip
   |-- Transient: dt=0.5us, T_total=100us (200 steps)
   |-- Output: P,V every 2 steps -> 100 frames
6. Solve -> 200 timesteps (12 cores, 41 min)
7. CFX-Post (Power Syntax)
   |-- Plane XZ at Y=0
   |-- Contour: Pressure, Velocity
   |-- Export: images (PNG) per timestep
   |-- ffmpeg -> animation.mp4
\end{lstlisting}

\subsection{Геометрiя (CadQuery $\to$ STEP)}

Скрипт \texttt{geometry/lab4.py}: ванночка (stadium/капсула, $220 \times 120 \times 150$~мм) з виглибком на днищi ($D=40$~мм, глибина~5~мм) --- поверхня випромiнювача. Одне рiдинне тiло.

\subsection{Матерiал: CompressibleWater}

\textbf{Критично}: бiблiотечна Water --- нестислива, не поширює звуковi хвилi. Створено власний матерiал з Linear Elastic EOS: $c = 1480$~м/с, $\rho = 998{,}2$~кг/м$^3$. Це забезпечує коректне моделювання поширення акустичних хвиль.

% ===================================================================
\section{Вiзуалiзацiя розрахункової моделi}
% ===================================================================

\begin{figure}[H]
  \centering
  \includegraphics[width=0.75\textwidth]{img/geometry}
  \caption{Геометрiя ультразвукової ванночки з випромiнювачем на днищi}
\end{figure}

На зображеннi видно ванночку капсулоподiбної форми з кiлькома скруглення на днищi та цилiндричним виглибком ($D=40$~мм) --- поверхню випромiнювача.

\begin{figure}[H]
  \centering
  \includegraphics[width=0.75\textwidth]{img/mesh}
  \caption{Розрахункова сiтка (Mesh), розмiр елемента 3~мм}
\end{figure}

Сiтка побудована з глобальним розмiром елемента 3~мм ($\lambda/21$), що забезпечує щонайменше 21 елемент на довжину хвилi --- значно бiльше мiнiмальних 10 елементiв. Загальна кiлькiсть: 200418 вузлiв, 1095393 тетраедральних елементи.

% ===================================================================
\section{Результати розрахунку}
% ===================================================================

\subsection{Pressure}

\begin{figure}[H]
  \centering
  \includegraphics[width=0.95\textwidth]{img/pressure}
  \caption{Contour: розподiл тиску (Pressure) на площинi XZ при $t = 100$~мкс}
\end{figure}

На контурi тиску чiтко видно концентричнi хвильовi фронти, що поширюються вiд випромiнювача (на днищi ванночки) вгору та в сторони. Максимальний тиск досягає 455~кПа у зонi стиснення безпосередньо над випромiнювачем. Чергування зон пiдвищеного тиску (стиснення) та фонового тиску (розрiдження) вiдповiдає хвильовому характеру процесу.

Вiдстань мiж сусiднiми хвильовими фронтами вiзуально вiдповiдає теоретичнiй довжинi хвилi $\lambda \approx 64$~мм.

\subsection{Velocity}

\begin{figure}[H]
  \centering
  \includegraphics[width=0.95\textwidth]{img/velocity}
  \caption{Contour: розподiл швидкостi (Velocity) на площинi XZ при $t = 100$~мкс}
\end{figure}

Розподiл швидкостi демонструє кiльцеву структуру, що корелює з полем тиску: максимальнi швидкостi (до 0{,}121~м/с) спостерiгаються безпосередньо над випромiнювачем, де градiєнт тиску найбiльший. Хвильовi фронти швидкостi також мають концентричну форму, що пiдтверджує сферичний характер поширення хвилi вiд точкового джерела.

\subsection{Кiлькiснi характеристики}

\begin{table}[H]
  \centering
  \caption{Параметри хвильового поля при $t = 100$~мкс}
  \begin{tabularx}{0.80\textwidth}{Xrr}
    \toprule
    \textbf{Параметр} & \textbf{Мiнiмум} & \textbf{Максимум} \\
    \midrule
    Тиск, кПа              & $-390$  & $+455$ \\
    Швидкiсть, м/с          & 0       & 0{,}121 \\
    Змiщення випромiнювача, нм & $-28{,}5$ & 0 \\
    Число Маха              & ---     & $1{,}5 \times 10^{-5}$ \\
    Число Рейнольдса        & ---     & 3852 \\
    \bottomrule
  \end{tabularx}
\end{table}

Максимальне змiщення випромiнювача складає 28{,}5~нм, що вiдповiдає задаонiй амплiтудi 1{,}3~мкм з урахуванням кiлькостi пройдених перiодiв. Число Маха $\ll 1$ пiдтверджує лiнiйний акустичний режим.

\subsection{Анiмацiя}

Створено вiдеофайл \texttt{CFX\_lab4.mp4}, що демонструє динамiку поширення ультразвукових хвиль протягом 200 часових крокiв (100~мкс).

% ===================================================================
\section{Аналiз та висновок}
% ===================================================================

\subsection{Хвильовi характеристики}

Аналiз результатiв дозволяє видiлити наступнi особливостi:

\begin{itemize}
  \item \textbf{Хвилi стиснення/розрiдження} --- чiтко вiзуалiзуються як чергування зон пiдвищеного та зниженого тиску. Амплiтуда тиску: вiд $-390$~кПа до $+455$~кПа вiдносно атмосферного.

  \item \textbf{Довжина хвилi} --- вiдстань мiж сусiднiми фронтами на контурi тиску вiзуально вiдповiдає теоретичному значенню $\lambda = c/f = 1480/23000 \approx 64{,}3$~мм. За 200 часових крокiв ($T = 100$~мкс) хвиля проходить вiдстань $c \cdot T = 1480 \cdot 10^{-4} = 0{,}148$~м $\approx 2{,}3\lambda$, що узгоджується з кiлькiстю видимих фронтiв.

  \item \textbf{Сферичне поширення} --- хвильовi фронти мають напiвсферичну форму, що вiдповiдає поширенню вiд малого плоского джерела ($D_{\mathrm{emitter}} = 40$~мм $< \lambda$). Оскiльки $D < \lambda$, випромiнювач дiє як точкове джерело.

  \item \textbf{Вiдбиття} --- при досягненнi стiнок ванночки хвилi вiдбиваються, формуючи iнтерференцiйну картину. На пiзнiших кадрах анiмацiї видно ускладнення хвильового поля.
\end{itemize}

\subsection{Параметри розрахунку}

\begin{itemize}
  \item Часовий крок $\Delta t = 0{,}5$~мкс забезпечує $\sim 87$ крокiв на перiод хвилi ($T = 1/f = 43{,}5$~мкс), що значно перевищує мiнiмальнi 20 крокiв.
  \item Розмiр елемента 3~мм ($\lambda/21$) забезпечує адекватне просторове розрiшення хвильового фронту.
  \item Акустичне число Куранта $C_a = c \cdot \Delta t / h = 1480 \cdot 5 \times 10^{-7} / 0{,}003 = 0{,}247$ --- менше 1, що забезпечує стабiльнiсть явної схеми.
  \item Загальний час розрахунку: 41~хв на 12 ядрах.
\end{itemize}

\textbf{Висновок:} Моделювання ультразвукової ванночки при $f = 23$~кГц, $A = 1{,}3$~мкм демонструє фiзично реалiстичну картину поширення акустичних хвиль. Чiтко вiзуалiзуються хвилi стиснення та розрiдження з довжиною хвилi $\lambda \approx 64$~мм, що вiдповiдає теоретичному значенню для води ($c = 1480$~м/с). Амплiтуда тиску досягає $\pm 400$~кПа, швидкiсть часток --- до 0{,}12~м/с. Ванночка об'ємом $\sim 3{,}6$~л ефективно заповнюється акустичним полем вже за $\sim 2$ перiоди хвилi ($\sim 87$~мкс). Використання CompressibleWater з Linear Elastic EOS є критичним для коректного моделювання --- стандартна нестислива Water не поширює звуковi хвилi.

% ===================================================================
\appendix
\section{Скрипт генерацiї геометрiї}
% ===================================================================

\begin{lstlisting}[style=python]
"""Lab 4 -- Ultrasonic bath geometry, Variant 13.

Stadium/capsule-shaped bath with:
  - Filleted bottom edges
  - Circular emitter recess (D=40mm, 5mm deep)
  - Emitter face = flat circle at Z=5mm
"""

from pathlib import Path
import cadquery as cq

BATH_WIDTH = 120.0
BATH_STRAIGHT = 100.0
BATH_LENGTH = BATH_STRAIGHT + BATH_WIDTH  # 220 mm
BATH_HEIGHT = 150.0
BOTTOM_FILLET = 10.0

EMITTER_D = 40.0
EMITTER_DEPTH = 5.0

bath = (
    cq.Workplane("XY")
    .slot2D(BATH_LENGTH, BATH_WIDTH)
    .extrude(BATH_HEIGHT)
)

bath = bath.edges("<Z").fillet(BOTTOM_FILLET)

emitter_cut = (
    cq.Workplane("XY")
    .circle(EMITTER_D / 2)
    .extrude(EMITTER_DEPTH)
)
bath = bath.cut(emitter_cut)

step_dir = Path(__file__).parent / "step"
step_dir.mkdir(exist_ok=True)
step_path = step_dir / "lab4.step"
cq.exporters.export(bath, str(step_path))
\end{lstlisting}

\end{document}
